% -----------------------------------------------------------------------------
%  Figure 3.5 — Chaîne d'intégration et de déploiement continus en service.
%  Bande du haut : l'enchaînement déclenché à chaque dépôt de modification.
%  Bande du bas : le déploiement, réservé à la branche principale, qui se termine
%  par la vérification du site réellement servi au visiteur.
% -----------------------------------------------------------------------------
\begin{tikzpicture}[x=1cm, y=1cm]

  % ===================== INTÉGRATION CONTINUE ================================
  \node[blocfort, text width=2.3cm, anchor=north, minimum height=13mm] (dev)
       at (-6.0,0) {Poste de développement\\ \texttt{git push}};

  \node[bloc, text width=2.2cm, anchor=north, minimum height=13mm] (git)
       at (-3.2,0) {\textbf{Dépôt GitHub}};

  \node[bloc, text width=3.1cm, anchor=north, minimum height=13mm] (ci)
       at (0.2,0) {\textbf{Intégration continue}\\
                   installation, compilation,\\
                   \texttt{terraform fmt} et \texttt{validate}};

  \node[bloc, text width=3.4cm, anchor=north, minimum height=13mm] (ctrl)
       at (4.0,0) {\textbf{Contrôles du paquet produit}\\
                   fichiers attendus, modèles 3D valides,
                   adresse du service de données};

  \draw[fleche] (dev.east) -- (git.west);
  \draw[fleche] (git.east) -- (ci.west);
  \draw[fleche] (ci.east) -- (ctrl.west);

  \begin{scope}[on background layer]
    \node[couche, fit=(dev)(ctrl), draw=keycemid] (bci) {};
  \end{scope}
  \node[titrecouche, anchor=west] at ([xshift=6pt]bci.north west)
       {À CHAQUE DÉPÔT DE MODIFICATION};

  % ===================== CONDITION D'ARRÊT ===================================
  \node[blocgris, text width=3.4cm, anchor=north, minimum height=9mm] (stop)
       at (4.0,-3.0) {\textbf{Échec} : la chaîne s'arrête,\\ rien n'est publié};
  \draw[fleche, draw=keyceorange] (ctrl.south) -- (stop.north)
        node[etiquette, pos=0.5, anchor=east, xshift=-3pt, text=keyceorange]
        {si un contrôle\\ échoue};

  % ===================== DÉPLOIEMENT CONTINU =================================
  \node[bloc, text width=2.5cm, anchor=north, minimum height=14mm] (oidc)
       at (-6.0,-4.9) {\textbf{Jeton OIDC}\\ de courte durée\\
                       \textit{aucune clé permanente}};

  \node[bloc, text width=2.4cm, anchor=north, minimum height=14mm] (pub)
       at (-2.9,-4.9) {Publication en quatre passes vers \textbf{S3}};

  \node[bloc, text width=2.4cm, anchor=north, minimum height=14mm] (purge)
       at (0.2,-4.9) {Purge du cache \textbf{CloudFront}};

  \node[blocfort, text width=3.2cm, anchor=north, minimum height=14mm] (verif)
       at (3.8,-4.9) {\textbf{Vérification du site en ligne}\\
                      \textit{on contrôle ce que le visiteur reçoit}};

  \draw[fleche] (oidc.east) -- (pub.west);
  \draw[fleche] (pub.east) -- (purge.west);
  \draw[fleche] (purge.east) -- (verif.west);

  \begin{scope}[on background layer]
    \node[couche, fit=(oidc)(verif), draw=keycemid] (bcd) {};
  \end{scope}
  \node[titrecouche, anchor=west] at ([xshift=6pt]bcd.north west)
       {SUR LA BRANCHE PRINCIPALE SEULEMENT};

  % --- passage d'une bande à l'autre -----------------------------------------
  \draw[fleche] ([xshift=1.1cm]bci.south west) -- ([xshift=1.1cm]bcd.north west)
        node[etiquette, pos=0.5, anchor=west, xshift=4pt]
        {si tous les contrôles passent};

  % ===================== VOIE DOCKER =========================================
  \node[blocext, text width=6.6cm, anchor=north, minimum height=10mm] (doc)
       at (-3.2,-7.4) {\textbf{Voie Docker} — image multi-étapes (compilation
                       Node, puis diffusion Nginx), employée pour reproduire en
                       local le comportement de production};
  \draw[flechefine, draw=keyceorange] (git.south) -- ++(0,-0.55) -| (doc.north);

\end{tikzpicture}
